\documentclass[12pt,a4paper,twoside]{report}

%% ─────────────────────────────────────────────────────────────────────────────
%% PACKAGES
%% ─────────────────────────────────────────────────────────────────────────────
\usepackage[T1]{fontenc}
\usepackage[utf8]{inputenc}
\usepackage{charter}
\usepackage{microtype}
\usepackage[margin=2.5cm, inner=3.0cm, outer=2.5cm]{geometry}

%% Mathematics
\usepackage{amsmath,amssymb,amsthm,mathtools}
\usepackage{physics}          % \dd, \pdv, \expval, \norm, etc.
\usepackage{bm}               % bold math

%% Units and numbers
\usepackage{siunitx}
\sisetup{separate-uncertainty=true, range-phrase=\text{--}}

%% Tables and figures
\usepackage{booktabs}
\usepackage{multirow}
\usepackage{array}
\usepackage{graphicx}
\usepackage{caption}
\usepackage{subcaption}
\usepackage{float}

%% Code listings
\usepackage{listings}
\usepackage{xcolor}

%% Algorithms
\usepackage[ruled,vlined,linesnumbered]{algorithm2e}

%% Hyperlinks and references
\usepackage[colorlinks=true,
            linkcolor=blue!60!black,
            citecolor=green!50!black,
            urlcolor=blue!70!black]{hyperref}
\usepackage{cleveref}

%% Miscellaneous
\usepackage{enumitem}
\usepackage{tcolorbox}
\tcbuselibrary{skins,breakable}
\usepackage{setspace}
\onehalfspacing

%% ─────────────────────────────────────────────────────────────────────────────
%% LISTING STYLE  (Python code blocks)
%% ─────────────────────────────────────────────────────────────────────────────
\definecolor{codebg}{HTML}{0d0f14}
\definecolor{codeframe}{HTML}{1e2230}
\definecolor{codecomment}{HTML}{6a9955}
\definecolor{codestring}{HTML}{ce9178}
\definecolor{codekw}{HTML}{4fc3f7}
\definecolor{codenumber}{HTML}{b5cea8}
\definecolor{codefunc}{HTML}{ffd54f}
\definecolor{codetext}{HTML}{d8dce8}

\lstdefinestyle{pythonstyle}{
  backgroundcolor=\color{codebg},
  basicstyle=\ttfamily\small\color{codetext},
  commentstyle=\color{codecomment}\itshape,
  stringstyle=\color{codestring},
  keywordstyle=\color{codekw}\bfseries,
  numberstyle=\tiny\color{codetext!50},
  breakatwhitespace=false,
  breaklines=true,
  captionpos=b,
  keepspaces=true,
  numbers=left,
  numbersep=8pt,
  showspaces=false,
  showstringspaces=false,
  showtabs=false,
  tabsize=4,
  frame=single,
  rulecolor=\color{codeframe},
  framesep=4pt,
  xleftmargin=14pt,
  xrightmargin=4pt,
  language=Python,
  morekeywords={True,False,None,yield,as,with,lambda,from,import},
  literate={*}{{\textcolor{codekw}{*}}}1
           {=}{{\textcolor{codekw}{=}}}1,
}
\lstset{style=pythonstyle}

%% ─────────────────────────────────────────────────────────────────────────────
%% CUSTOM COMMANDS
%% ─────────────────────────────────────────────────────────────────────────────
\newcommand{\trento}{\textsc{TRENTo}\xspace}
% \epn removed: conflict with physics package
\newcommand{\ep}[1]{\varepsilon_{#1}}
\newcommand{\epn}[2]{\varepsilon_{#1}\{#2\}}
\newcommand{\vn}[2]{v_{#1}\{#2\}}
\newcommand{\mean}[1]{\langle #1 \rangle}
\newcommand{\PbPb}{$^{208}$Pb$+^{208}$Pb\xspace}
\newcommand{\sqrtsnn}{\sqrt{s_\mathrm{NN}}}
\newcommand{\dNdeta}{dN_\mathrm{ch}/d\eta}
\newcommand{\WS}{\text{WS}}

%% Theorem-like environments
\newtheorem{definition}{Definition}[section]
\newtheorem{remark}{Remark}[section]

%% Coloured note boxes
\tcbset{
  physbox/.style={
    colback=blue!5!white, colframe=blue!60!black,
    fonttitle=\bfseries, title=#1,
    breakable, enhanced,
  },
  puzzlebox/.style={
    colback=red!5!white, colframe=red!60!black,
    fonttitle=\bfseries, title=#1,
    breakable, enhanced,
  },
  codebox/.style={
    colback=codebg, colframe=codeframe,
    fonttitle=\bfseries\color{codefunc}, title=#1,
    breakable, enhanced,
  },
}

%% ─────────────────────────────────────────────────────────────────────────────
%% DOCUMENT
%% ─────────────────────────────────────────────────────────────────────────────
\begin{document}

%%  Stand-alone compilation: fake chapter number
\setcounter{chapter}{3}

% ══════════════════════════════════════════════════════════════════════════════
\chapter{Initial-State Geometry and the Ultracentral Flow Puzzle in
         \texorpdfstring{\PbPb}{208Pb+208Pb} Collisions}
\label{ch:ultracentral}
% ══════════════════════════════════════════════════════════════════════════════

\begin{quote}
\itshape
This chapter describes the computational framework developed to investigate
the ultracentral \PbPb flow puzzle through a systematic scan of the nuclear
initial-state parameter space.  We introduce the theoretical foundations of
initial-state geometry, the nuclear structure model, the \trento
entropy-deposition framework, and the statistical design methodology.
The full pipeline — from Latin Hypercube Sampling through eccentricity
computation — is described in both mathematical and algorithmic terms.
\end{quote}

% ──────────────────────────────────────────────────────────────────────────────
\section{Introduction and Motivation}
\label{sec:intro_puzzle}
% ──────────────────────────────────────────────────────────────────────────────

The collective flow of matter produced in ultrarelativistic heavy-ion
collisions at the Large Hadron Collider (LHC) is among the most precisely
studied emergent phenomena in nuclear physics.  In non-central collisions
($b \gtrsim 2\,\si{fm}$), the initial almond-shaped overlap region of the
two colliding nuclei generates a dominant elliptic anisotropy in momentum
space, characterised by $\vn{2}{2}$.  As the collision centrality approaches
zero, the geometric contribution of the impact parameter vanishes, and all
flow harmonics are driven solely by quantum fluctuations in nucleon positions
within the individual nuclei.

In this regime — hereafter referred to as the \emph{ultracentral} class,
operationally defined as the \mbox{$0$--$1\%$} most central collisions — a
cluster of persistent discrepancies between experimental data and state-of-the-art
hydrodynamic model calculations has been identified.  Collectively, these
discrepancies constitute the \textbf{ultracentral flow puzzle}:

\begin{tcolorbox}[puzzlebox={The Ultracentral Flow Puzzle — Experimental Signatures}]
\begin{enumerate}[label=(\roman*), leftmargin=*]
  \item \textbf{$v_2$-to-$v_3$ hierarchy reversal.}
        Viscous hydrodynamic calculations universally predict
        $\vn{2}{2} > \vn{3}{2}$ in the $0$--$1\%$ centrality bin, because
        the elliptic eccentricity $\ep{2}$ retains a small finite mean from
        the residual geometry.  ALICE and ATLAS measurements instead show
        $\vn{2}{2} \approx \vn{3}{2}$, implying an unexpected suppression of
        elliptic flow or enhancement of triangular flow
        \cite{Aamodt2011,Aad2012,Luzum2013}.
  \item \textbf{Negative $v_4\{4\}^4$.}
        The four-particle hexadecapole cumulant combination
        $v_4\{4\}^4 \equiv 2\mean{v_4^2}^2 - \mean{v_4^4}$ is measured
        to be \emph{negative} in \PbPb at $\sqrtsnn = 2.76$ and
        $5.02\,\si{TeV}$ by ALICE, ATLAS, and CMS
        \cite{Adam2016,Aaboud2019,Sirunyan2017}.
        Standard models predict a positive value.
  \item \textbf{Small $v_2\{4\}/v_2\{2\}$ ratio.}
        The fluctuation-suppression ratio
        $R_{42} \equiv v_2\{4\}/v_2\{2\}$
        approaches values below $0.5$ in the $0$--$0.1\%$ bin,
        significantly smaller than model predictions, indicating
        a non-Gaussian distribution $p(\varepsilon_2)$ with much lighter
        tails than expected \cite{Giacalone2017}.
  \item \textbf{Enhanced $v_3\{4\}/v_3\{2\}$ ratio.}
        Models systematically underpredict this ratio in ultra-central
        collisions, i.e.\ they overestimate the relative width of
        $p(\varepsilon_3)$ \cite{Chatrchyan2014}.
\end{enumerate}
\end{tcolorbox}

\noindent
Proposed resolutions involve modifications to the nuclear structure of
$^{208}$Pb: a small octupole deformation $\beta_3$
\cite{Carzon2020}, a hexadecapole component $\beta_4$ coupled to the
nonlinear mode-mixing channel $v_4 \sim \chi_{42}\,v_2^2$
\cite{Giacalone2017b}, short-range nucleon-nucleon correlations
\cite{Luzum2023}, and modifications to the nucleon substructure
(Trento parameters $w$, $k$).  The primary goal of the present study
is to perform a systematic, quantitative scan of this parameter space
using a Latin Hypercube design, enabling subsequent Gaussian Process
emulation and Bayesian inference.

% ──────────────────────────────────────────────────────────────────────────────
\section{Nuclear Structure: The Deformed Woods-Saxon Model}
\label{sec:nuclear_structure}
% ──────────────────────────────────────────────────────────────────────────────

\subsection{The Woods-Saxon Density Profile}
\label{subsec:WS}

The nuclear one-body density of a spherically symmetric nucleus of mass
number $A$ is well described by the Fermi distribution, commonly known
as the Woods-Saxon (WS) profile \cite{WoodsSaxon1954}:
\begin{equation}
  \rho_0(r) = \frac{\rho_\mathrm{sat}}{1 + \exp\!\left(\dfrac{r - R}{a}\right)},
  \label{eq:WS_spherical}
\end{equation}
where $\rho_\mathrm{sat} \approx 0.16\,\si{fm^{-3}}$ is the nuclear
saturation density, $R$ is the half-density radius, and $a$ is the
diffuseness parameter governing the surface-to-interior transition.
For $^{208}$Pb the empirical values extracted from electron-scattering
experiments are
$R = 6.624\,\si{fm}$ and $a = 0.549\,\si{fm}$
\cite{deVries1987}.

\subsection{Deformed Density Profiles}
\label{subsec:deformed_WS}

For a deformed nucleus, the surface is described by a multipole
expansion in real spherical harmonics $Y_\ell^0(\theta, \phi)$.  The
deformed WS radius is
\begin{equation}
  R(\theta, \phi) = R \left[
    1 + \sum_{\ell=2}^{\infty} \beta_\ell\, Y_\ell^0(\theta, \phi)
  \right],
  \label{eq:WS_deformed_radius}
\end{equation}
where $\beta_\ell$ is the dimensionless deformation parameter of
multipolarity $\ell$, and the sum runs over even and odd multipoles
as permitted by nuclear symmetry.  Substituting
\cref{eq:WS_deformed_radius} into \cref{eq:WS_spherical} gives the
deformed density:
\begin{equation}
  \rho(\bm{r}) = \frac{\rho_\mathrm{sat}}{1 + \exp\!\left(
    \dfrac{|\bm{r}| - R(\theta,\phi)}{a}
  \right)}.
  \label{eq:WS_deformed_full}
\end{equation}

The deformation parameters relevant to this study are:
\begin{itemize}[leftmargin=*]
  \item \textbf{Quadrupole $\beta_2$:}  The leading-order prolate/oblate
        deformation.  For doubly-magic $^{208}$Pb, the ground-state
        quadrupole moment is negligible: $\beta_2 \approx 0$
        \cite{Roca-Maza2018}.
  \item \textbf{Octupole $\beta_3$:}  A pear-shaped deformation that breaks
        reflection symmetry.  Although $^{208}$Pb is doubly magic, low-energy
        nuclear structure calculations permit a small dynamic octupole
        component.  Data from ALICE and ATLAS constrain
        $\beta_3 \lesssim 0.0375$ for static deformation \cite{Carzon2020};
        recent dynamical proposals allow an effective value up to
        $\beta_3 \sim 0.07$ \cite{Xu2025}.
  \item \textbf{Hexadecapole $\beta_4$:}  Couples to the fourth-order flow
        coefficient via the nonlinear channel.  DFT calculations give
        $|\beta_4| \sim 0.01$--$0.02$ for $^{208}$Pb \cite{Bally2022}.
\end{itemize}

For a nucleus undergoing a random three-dimensional rotation before
each collision (as appropriate for a quantum ground state), the
deformation enters the initial-state geometry only through the
event-by-event orientation average.  In the isobar-sampler framework
\cite{Luzum2023_sampler}, an Euler rotation is applied to each
generated nuclear configuration.

\subsection{Short-Range Nucleon-Nucleon Correlations}
\label{subsec:correlations}

Beyond the mean-field WS profile, nucleons interact via a hard-core
short-range repulsion at distances $r \lesssim 0.4\,\si{fm}$, arising
from Pauli blocking and the repulsive core of the nucleon-nucleon
potential.  This suppresses nucleon clustering at short distances and
modifies the sub-nucleonic granularity of the initial state.

The correlation function is implemented in the isobar sampler as a
step function:
\begin{equation}
  C(r_{ij}) =
  \begin{cases}
    C_s & \text{if } r_{ij} < C_\ell, \\
    0   & \text{if } r_{ij} \geq C_\ell,
  \end{cases}
  \label{eq:corr_function}
\end{equation}
where $r_{ij}$ is the separation between nucleons $i$ and $j$, $C_\ell$
is the correlation length (hard-core radius), and $C_s \in [-1, 0]$ is
the correlation strength.  The value $C_s = -1$ corresponds to full
Pauli exclusion (no two nucleons within $C_\ell$ of each other);
$C_s = 0$ is the uncorrelated Woods-Saxon case.  Following
\cite{Luzum2023_sampler}, we fix $C_\ell = 0.4\,\si{fm}$ and
$C_s = -1$ as the default values.

\subsection{The Isobar Sampler}
\label{subsec:isobar_sampler}

Generating nuclear configurations with varied WS parameters
independently for each design point would introduce uncorrelated
statistical noise that obscures the genuine physical response to
parameter changes.  The isobar-sampler framework of
Luzum~et~al.~\cite{Luzum2023_sampler} addresses this by factorising
the stochastic content into a \emph{seed bank}.

\paragraph{Stage 1: Seed generation.}
A bank of $N_\mathrm{conf}$ \emph{seeds} — vectors of independent
random numbers in $[0,1]$ — is generated once and stored.  These seeds
encode the relative positions of nucleons in a parameter-independent
space.  Formally, nucleon position $\bm{r}_i$ is computed as a
deterministic function of the WS parameters and the corresponding seed
$\xi_i$:
\begin{equation}
  \bm{r}_i = \mathcal{F}\!\left(R, a, \beta_2, \beta_3, \beta_4,
              C_\ell, C_s;\; \xi_i\right).
  \label{eq:seed_map}
\end{equation}

\paragraph{Stage 2: Configuration generation.}
For each design point $\theta = (R, a, \beta_3, \beta_4, \ldots)$,
the same seed bank is used with the new parameters to produce
$N_\mathrm{conf}$ nuclear configurations.  Because the mapping
\cref{eq:seed_map} is monotone in each seed component,
configurations generated from the same seed bank with different
parameters differ only through their physical response and not through
spurious statistical fluctuations.  The $N_\mathrm{conf}$ configurations
are written to an HDF5 file (e.g.\ \texttt{WS1.hdf}) for use by \trento.

% ──────────────────────────────────────────────────────────────────────────────
\section{Initial-State Entropy Deposition: The \trento Model}
\label{sec:trento}
% ──────────────────────────────────────────────────────────────────────────────

\subsection{General Entropy-Deposition Ansatz}
\label{subsec:trento_ansatz}

The \trento (Reduced Thickness Event-by-event Nuclear Topology) model
\cite{Moreland2015} provides a flexible parametric ansatz for the
initial-state entropy density $s(\bm{x}_\perp)$ in the transverse
plane.  The fundamental object is the \emph{reduced thickness function}:
\begin{equation}
  T_R(p;\, T_A, T_B)
  = \left(\frac{T_A^p + T_B^p}{2}\right)^{1/p},
  \label{eq:TR_general}
\end{equation}
where $T_A(\bm{x}_\perp)$ and $T_B(\bm{x}_\perp)$ are the participant
thickness functions of the two nuclei evaluated at transverse position
$\bm{x}_\perp$, and $p$ is the entropy-deposition exponent.
The entropy density is then
\begin{equation}
  s(\bm{x}_\perp) = \kappa\, T_R(p;\, T_A, T_B),
  \label{eq:trento_entropy}
\end{equation}
with $\kappa$ an overall normalisation constant.

The parameter $p$ interpolates between well-known limiting cases:
\begin{align}
  p \to -\infty &:\quad T_R \propto \min(T_A, T_B) \quad \text{(wounded nucleon, target)}, \\
  p = -1         &:\quad T_R \propto \frac{2\, T_A T_B}{T_A + T_B} \quad \text{(harmonic mean)}, \\
  p = 0          &:\quad T_R \propto \sqrt{T_A T_B} \quad \text{(geometric mean, KLN-like)}, \\
  p = +1         &:\quad T_R \propto \frac{T_A + T_B}{2} \quad \text{(arithmetic mean, WN-like)}.
  \label{eq:trento_limits}
\end{align}
Bayesian analyses of LHC Pb+Pb data prefer values near $p \approx 0$
\cite{Bernhard2019,Nijs2021}.

\subsection{Nucleon Sub-Structure}
\label{subsec:trento_nucleon}

The participant thickness function of nucleus $A$ is the sum over
participant nucleon contributions:
\begin{equation}
  T_A(\bm{x}_\perp) = \sum_{i \in \mathrm{part.}}
    \gamma_i\, \rho_\mathrm{nuc}(\bm{x}_\perp - \bm{s}_i),
  \label{eq:TA}
\end{equation}
where $\bm{s}_i$ is the transverse position of nucleon $i$,
$\gamma_i$ is an independent gamma-distributed fluctuation,
\begin{equation}
  \gamma_i \sim \mathrm{Gamma}(k, k^{-1}),
  \quad \mean{\gamma_i} = 1,\quad \mathrm{Var}(\gamma_i) = k^{-1},
  \label{eq:gamma_fluctuation}
\end{equation}
and $\rho_\mathrm{nuc}$ is the nucleon density profile.  Each nucleon
is modelled as a Gaussian:
\begin{equation}
  \rho_\mathrm{nuc}(\bm{x}_\perp) = \frac{1}{2\pi w^2}
  \exp\!\left(-\frac{|\bm{x}_\perp|^2}{2w^2}\right),
  \label{eq:nucleon_profile}
\end{equation}
where $w$ is the nucleon width in the transverse plane.  Smaller $w$
produces a grainier, more fluctuation-dominated initial state.

Two nucleons $i$ (in nucleus $A$) and $j$ (in nucleus $B$) are
considered to have interacted (participated) if their transverse
separation satisfies
\begin{equation}
  |\bm{s}_i - \bm{s}_j|^2 < d^2,
  \label{eq:collision_criterion}
\end{equation}
where $d$ is the minimum nucleon distance parameter.

\subsection{Fixed Trento Parameters}
\label{subsec:fixed_trento}

The Bayesian analyses of \cite{Bernhard2019} and \cite{Nijs2021}
provide well-constrained posterior distributions for the Trento
parameters in Pb+Pb.  In this study we fix:
\begin{equation}
  p = 0.0, \quad k = 1.6, \quad d = 1.0\,\si{fm},
  \quad \mathrm{norm} =
  \begin{cases}
    18.0 & \sqrtsnn = 2.76\,\si{TeV}, \\
    20.0 & \sqrtsnn = 5.02\,\si{TeV}.
  \end{cases}
  \label{eq:fixed_trento}
\end{equation}
The energy-dependent normalisation reflects the approximately
$\SI{11}{\percent}$ increase in charged-particle pseudorapidity density
$\dNdeta$ between the two LHC energies
\cite{ALICE_2.76,ALICE_5.02}.
Only the nucleon width $w$ is kept free, as it directly controls the
sub-nucleonic granularity of the initial state and hence the
fluctuation spectrum of the eccentricities.

% ──────────────────────────────────────────────────────────────────────────────
\section{Eccentricity Cumulants and the Flow Observables}
\label{sec:eccentricities}
% ──────────────────────────────────────────────────────────────────────────────

\subsection{Spatial Eccentricity Vector}
\label{subsec:eccentricity_def}

The spatial anisotropy of the initial-state entropy density is
characterised by the complex eccentricity vector of order $n$:
\begin{equation}
  \mathcal{E}_n = \varepsilon_n\, e^{in\Psi_n}
  \equiv -\frac{\mean{r^n\, e^{in\phi}}_s}{\mean{r^n}_s},
  \label{eq:eccentricity_vector}
\end{equation}
where the bracket $\mean{\cdot}_s$ denotes the entropy-weighted
two-dimensional average over the transverse plane:
\begin{equation}
  \mean{f}_s = \frac{\int f(\bm{x}_\perp)\, s(\bm{x}_\perp)\,
    d^2\bm{x}_\perp}{\int s(\bm{x}_\perp)\,d^2\bm{x}_\perp},
  \label{eq:entropy_weight}
\end{equation}
$(r, \phi)$ are polar coordinates in the participant-centred frame,
$\varepsilon_n = |\mathcal{E}_n| \in [0,1]$ is the eccentricity
magnitude, and $\Psi_n$ is the $n$-th order participant-plane angle.
The $r^n$ weighting (Teaney--Yan prescription \cite{Teaney2011})
ensures a linear response between $\varepsilon_n$ and $v_n$
for $n = 2, 3$ at leading order in viscous corrections.

\subsection{Multi-Particle Cumulants of the Eccentricity Distribution}
\label{subsec:eccentricity_cumulants}

Let $\varepsilon_n$ denote the magnitude of the eccentricity in a given
event.  Across an ensemble of events within a centrality class, the
distribution $p(\varepsilon_n)$ encodes both the mean geometry and its
fluctuations.  Multi-particle flow cumulants measure moments of this
distribution in a way that is robust against non-flow correlations
\cite{Borghini2001}.  The eccentricity cumulants are defined through
the moment generating function:
\begin{align}
  \varepsilon_n\{2\} &= \sqrt{\mean{\varepsilon_n^2}},
    \label{eq:e2_cumulant} \\[4pt]
  \varepsilon_n\{4\} &= \left(2\mean{\varepsilon_n^2}^2 -
    \mean{\varepsilon_n^4}\right)^{1/4},
    \label{eq:e4_cumulant} \\[4pt]
  \varepsilon_n\{6\} &= \left(-\frac{c_6}{4}\right)^{1/6},\quad
    c_6 = \mean{\varepsilon_n^6} - 9\mean{\varepsilon_n^4}\mean{\varepsilon_n^2}
          + 12\mean{\varepsilon_n^2}^3,
    \label{eq:e6_cumulant} \\[4pt]
  \varepsilon_n\{8\} &= \left(\frac{c_8}{33}\right)^{1/8},\quad
    c_8 = \mean{\varepsilon_n^8} - 16\mean{\varepsilon_n^6}\mean{\varepsilon_n^2}
          - 18\mean{\varepsilon_n^4}^2
          + 144\mean{\varepsilon_n^4}\mean{\varepsilon_n^2}^2
          - 144\mean{\varepsilon_n^2}^4.
    \label{eq:e8_cumulant}
\end{align}
The convergence $\varepsilon_n\{4\} = \varepsilon_n\{6\} = \varepsilon_n\{8\}$
signals a Gaussian distribution $p(\varepsilon_n)$.  Deviations indicate
non-Gaussianity, which is the hallmark of the ultracentral regime.

\subsection{Puzzle-Critical Observables}
\label{subsec:puzzle_observables}

We now define the four observables that directly diagnose the
ultracentral flow puzzle, together with their eccentricity proxies.

\subsubsection{Triangularity Ratio \texorpdfstring{$T_{32}$}{T32}}

The ratio
\begin{equation}
  T_{32} \equiv \frac{\varepsilon_3\{2\}}{\varepsilon_2\{2\}}
  = \frac{\sqrt{\mean{\varepsilon_3^2}}}{\sqrt{\mean{\varepsilon_2^2}}}
  \label{eq:T32}
\end{equation}
tracks the transition from the geometry-dominated regime
($T_{32} \ll 1$ in peripheral collisions, where $\varepsilon_2$ is large)
to the fluctuation-dominated ultracentral regime.  Spherical models predict
$T_{32} \lesssim 0.8$ in the $0$--$1\%$ bin; ALICE data are consistent with
$T_{32} \approx 1.0$ (see Sec.~\ref{sec:intro_puzzle}).

\subsubsection{Fluctuation-Suppression Ratio \texorpdfstring{$R_{42}$}{R42}}

\begin{equation}
  R_{42} \equiv \frac{\varepsilon_2\{4\}}{\varepsilon_2\{2\}}
  = \left(\frac{2\mean{\varepsilon_2^2}^2 - \mean{\varepsilon_2^4}}
               {\mean{\varepsilon_2^2}^2}\right)^{1/4}
  \label{eq:R42}
\end{equation}
is bounded $R_{42} \leq 1$ and equals unity in the absence of
fluctuations.  For a distribution with mean $\bar{\varepsilon}_2$
and variance $\sigma^2$:
\begin{equation}
  R_{42} \approx \frac{\bar{\varepsilon}_2}
                       {\sqrt{\bar{\varepsilon}_2^2 + \sigma^2}}.
  \label{eq:R42_approx}
\end{equation}
As $\bar{\varepsilon}_2 \to 0$ in ultra-central collisions,
$R_{42} \to 0$ unless $\sigma^2$ also vanishes.  Data in the $0$--$0.1\%$
bin show $R_{42} < 0.5$, indicating a strongly non-Gaussian
$p(\varepsilon_2)$ with lighter tails than any current model predicts.

\subsubsection{Normalised Kurtosis \texorpdfstring{$\Gamma_2$}{Gamma2}}

The normalised excess kurtosis of $p(\varepsilon_2)$,
\begin{equation}
  \Gamma_2 \equiv
  \frac{\mean{\varepsilon_2^4} - 2\mean{\varepsilon_2^2}^2}
       {\mean{\varepsilon_2^2}^2},
  \label{eq:Gamma2}
\end{equation}
satisfies $\Gamma_2 = 0$ for a Gaussian,
$\Gamma_2 > 0$ for heavier tails (leptokurtic), and
$\Gamma_2 < 0$ for lighter tails (platykurtic).
In the ultracentral limit where $p(\varepsilon_2)$ approaches a
Rayleigh distribution, $\Gamma_2 \to -2$, i.e.\
$\mean{\varepsilon_2^4} \to 2\mean{\varepsilon_2^2}^2 / (1+1)
= \mean{\varepsilon_2^2}^2$.
The relation to the fluctuation-suppression ratio is:
\begin{equation}
  R_{42}^4 = 2 + \Gamma_2.
  \label{eq:R42_Gamma2}
\end{equation}

\subsubsection{Hexadecapole Sign-Change Observable
  \texorpdfstring{$\varepsilon_4\{4\}^4$}{e4(4)4}}

The four-particle hexadecapole cumulant combination:
\begin{equation}
  \varepsilon_4\{4\}^4 \equiv 2\mean{\varepsilon_4^2}^2 - \mean{\varepsilon_4^4}
  \label{eq:e4_4c}
\end{equation}
is the direct eccentricity analogue of the measured $v_4\{4\}^4$.
This quantity becomes \emph{negative} when the distribution $p(\varepsilon_4)$
is sufficiently broad, i.e.\ when $\mean{\varepsilon_4^4} > 2\mean{\varepsilon_4^2}^2$.

The mechanism driving $v_4\{4\}^4 < 0$ involves the nonlinear
mode-coupling channel $v_4 \approx \chi_{42}\,v_2^2$, where
$\chi_{42}$ is the linear response coefficient.  The contribution of
this channel to $v_4\{4\}^4$ is negative whenever
\cite{Giacalone2017b}:
\begin{equation}
  \frac{\mean{\varepsilon_2^8}}{\mean{\varepsilon_2^4}^2} > 2.
  \label{eq:sign_change_trigger}
\end{equation}
This ratio — denoted \texttt{nl\_r8} in the code — therefore serves as
a \emph{sign-change trigger}: if it exceeds $2$, the predicted
$v_4\{4\}^4$ will be negative even for a model with vanishing mean
$\varepsilon_4$.  The nonlinear eccentricity coupling is:
\begin{equation}
  \varepsilon_4^{\Psi_2} \equiv
  \frac{\mean{\varepsilon_4\, \varepsilon_2^2}}{\sqrt{\mean{\varepsilon_2^4}}},
  \label{eq:e4_psi2}
\end{equation}
which measures the correlation between the hexadecapole and elliptic
eccentricities in the same event.

\subsection{Summary of Computed Observables}
\label{subsec:obs_summary}

\Cref{tab:observables} summarises the complete set of eccentricity
observables computed for each design point in the $0$--$1\%$
ultracentral bin.

\begin{table}[htbp]
\centering
\caption{Eccentricity observables computed per design point in the
  $0$--$1\%$ centrality bin.  All are evaluated using the event ensemble
  from a single Trento run.  ``KEY'' denotes a puzzle-critical quantity.}
\label{tab:observables}
\begin{tabular}{llll}
\toprule
Code key & Mathematical definition & Physical meaning & Puzzle? \\
\midrule
\texttt{e2\_2} & $\sqrt{\mean{\varepsilon_2^2}}$ & Elliptic eccentricity (RMS) & \\
\texttt{e2\_4} & $(2\mean{\varepsilon_2^2}^2 - \mean{\varepsilon_2^4})^{1/4}$ & 4-particle cumulant & \\
\texttt{e2\_6} & $(-c_6/4)^{1/6}$ & 6-particle cumulant & \\
\texttt{e2\_8} & $(c_8/33)^{1/8}$ & 8-particle cumulant & \\
\texttt{e3\_2} & $\sqrt{\mean{\varepsilon_3^2}}$ & Triangularity (RMS) & KEY \\
\texttt{e4\_2} & $\sqrt{\mean{\varepsilon_4^2}}$ & Hexadecapole (RMS) & \\
\texttt{R42}   & $\varepsilon_2\{4\}/\varepsilon_2\{2\}$ & Fluctuation suppression & KEY \\
\texttt{T32}   & $\varepsilon_3\{2\}/\varepsilon_2\{2\}$ & Triangularity ratio & KEY \\
\texttt{Gamma2}& $(\mean{\varepsilon_2^4} - 2\mean{\varepsilon_2^2}^2)/\mean{\varepsilon_2^2}^2$ & Excess kurtosis & KEY \\
\texttt{e4\_4c}& $2\mean{\varepsilon_4^2}^2 - \mean{\varepsilon_4^4}$ & Hexadecapole sign-change & KEY \\
\texttt{nl\_r8}& $\mean{\varepsilon_2^8}/\mean{\varepsilon_2^4}^2$ & NL sign-change trigger & KEY \\
\bottomrule
\end{tabular}
\end{table}

% ──────────────────────────────────────────────────────────────────────────────
\section{Centrality Determination}
\label{sec:centrality}
% ──────────────────────────────────────────────────────────────────────────────

Centrality is determined from the charged-particle multiplicity
distribution $dN/d\eta$ via a mapping from the Glauber model.
In the \trento simulation, each event produces a multiplicity proxy
proportional to $\int s(\bm{x}_\perp)\,d^2\bm{x}_\perp$.  Events are
ranked in descending order of multiplicity $M$, and the centrality
percentile of event $i$ is assigned as:
\begin{equation}
  c_i = \frac{\mathrm{rank}(-M_i)}{N_\mathrm{ev}} \times 100\%,
  \label{eq:centrality}
\end{equation}
where $\mathrm{rank}(-M_i)$ is the rank of event $i$ in the
descending multiplicity ordering and $N_\mathrm{ev}$ is the total
number of events.  The ultracentral bin is defined as
$c \in [0\%, 1\%]$, and the bin population for $N_\mathrm{ev} = 5 \times 10^5$
is $N_\mathrm{UC} \approx 5000$ events, sufficient for a statistical
uncertainty of $\lesssim 1\%$ on all cumulant observables after
bootstrap resampling.

% ──────────────────────────────────────────────────────────────────────────────
\section{Statistical Design: The Latin Hypercube}
\label{sec:lhs}
% ──────────────────────────────────────────────────────────────────────────────

\subsection{The Design Problem}
\label{subsec:design_problem}

We wish to evaluate the map
\begin{equation}
  \bm{\theta} = (R,\, a,\, \beta_3,\, \beta_4,\, w)
  \longrightarrow
  \bm{y}(\bm{\theta}) = \bigl(\varepsilon_2\{2\},\,
    T_{32},\, R_{42},\, \Gamma_2,\, \varepsilon_4\{4\}^4, \ldots\bigr)
  \label{eq:parameter_to_observable_map}
\end{equation}
at $N_\mathrm{d}$ design points $\{\bm{\theta}^{(1)},
\ldots, \bm{\theta}^{(N_\mathrm{d})}\}$ distributed across the
five-dimensional parameter space:
\begin{equation}
  \bm{\theta} \in \Theta \equiv [6.50, 6.80]
    \times [0.44, 0.65]
    \times [0.00, 0.12]
    \times [-0.02, 0.06]
    \times [0.40, 1.20]
    \quad \text{(fm, fm, —, —, fm)}.
  \label{eq:prior_box}
\end{equation}
These ranges are motivated by the nuclear data compilations of
\cite{deVries1987,Carzon2020,Bally2022} and the Trento Bayesian
analysis of \cite{Bernhard2019}.

\subsection{Latin Hypercube Sampling}
\label{subsec:lhs_definition}

\begin{definition}[Latin Hypercube Sample]
A \emph{Latin Hypercube Sample} (LHS) of size $N$ in $d$ dimensions
is a matrix $\mathbf{U} \in [0,1]^{N \times d}$ such that, for each
column $j = 1, \ldots, d$, the entries
$\{U_{1j}, \ldots, U_{Nj}\}$ form a \emph{stratified} sample: each
stratum $[(k-1)/N,\, k/N)$ for $k = 1, \ldots, N$ contains exactly
one entry.
\end{definition}

A LHS achieves uniformity in every one-dimensional projection, which
is necessary but not sufficient for good coverage of the full
$d$-dimensional space.  To improve space-filling, we employ the
\emph{maximin} criterion \cite{Morris1995}:
\begin{equation}
  \mathbf{U}^* = \operatorname*{argmax}_{\mathbf{U} \in \mathcal{L}_N^d}
    \min_{i \neq j} \|\bm{u}_i - \bm{u}_j\|_2,
  \label{eq:maximin}
\end{equation}
where $\mathcal{L}_N^d$ is the set of all $N$-point LHS designs in
$d$ dimensions, and $\bm{u}_i$ is the $i$-th row of $\mathbf{U}$.

Because an exact solution of \cref{eq:maximin} is NP-hard for large
$N$ and $d$, we use a random search approximation:
$N_\mathrm{iter} = 2000$ independent random LHS candidates are
generated and the one with the largest minimum pairwise distance
is retained.

\paragraph{Design matrix.}
The physical design is obtained by affine scaling:
\begin{equation}
  \theta_j^{(i)} = \theta_j^\mathrm{min}
    + U_{ij}\, \bigl(\theta_j^\mathrm{max} - \theta_j^\mathrm{min}\bigr),
  \quad i = 1,\ldots,N;\; j = 1,\ldots,d.
  \label{eq:lhs_scaling}
\end{equation}

\subsection{Number of Design Points}
\label{subsec:n_design}

For a subsequent Gaussian Process (GP) emulator of the map
\cref{eq:parameter_to_observable_map} in $d = 5$ dimensions, the
standard recommendation is $N_\mathrm{d} \geq 10\,d = 50$ design
points as an absolute minimum, with $N_\mathrm{d} = 100$ providing
adequate coverage for a squared-exponential kernel with $\mathcal{O}(d)$
hyper-parameters \cite{Loeppky2009}.  We choose $N_\mathrm{d} = 100$
as our baseline and note that $N_\mathrm{d} = 200$ would be preferable
if the computational budget allows, as it better resolves non-linear
structures in parameter space.

\begin{remark}
The computational cost of the full pipeline scales as
$N_\mathrm{d} \times t_\mathrm{Trento}$, where
$t_\mathrm{Trento} \approx 15$--$30\,\si{min}$ for $5 \times 10^5$
events on 32 CPU cores.  For $N_\mathrm{d} = 100$, the total wall time
is approximately $25$--$50\,\si{hours}$; this is the primary
computational bottleneck.
\end{remark}

\subsection{Space-Filling Quality Assessment}
\label{subsec:lhs_quality}

The quality of the LHS is assessed by three complementary metrics:
\begin{enumerate}[label=(\alph*), leftmargin=*]
  \item \textbf{Marginal uniformity:} Histogram of each parameter's
        marginal distribution should be approximately flat.
  \item \textbf{Minimum pairwise distance:}
        $d_\mathrm{min} = \min_{i\neq j}\|\bm{u}_i - \bm{u}_j\|_2$
        in the unit hypercube; larger is better.
  \item \textbf{2-D projections:} The lower triangle of the corner
        plot (scatter of all $\binom{d}{2}$ parameter pairs) should
        show no obvious clustering or regular structure.
\end{enumerate}

\noindent
The corner plot produced by the pipeline (\texttt{corner\_lhs.pdf}) is
the primary visual diagnostic.

% ──────────────────────────────────────────────────────────────────────────────
\section{The Computational Pipeline}
\label{sec:pipeline}
% ──────────────────────────────────────────────────────────────────────────────

\subsection{Pipeline Architecture}
\label{subsec:pipeline_arch}

The full analysis is implemented in the Python script
\texttt{pbpb\_ultracentral\_scan.py}, which executes five sequentially
dependent stages summarised in \cref{tab:pipeline_stages}.  Each stage
writes its output to a persistent cache so that any individual stage
can be re-run independently by setting the corresponding
\texttt{RUN\_STAGE\_*} flag to \texttt{True}.

\begin{table}[htbp]
\centering
\caption{Pipeline stages, inputs, and outputs.}
\label{tab:pipeline_stages}
\begin{tabular}{p{0.4cm} p{3.2cm} p{3.6cm} p{4.2cm}}
\toprule
\# & Stage & Input & Output \\
\midrule
0 & LHS design generation &
  Configuration flags &
  \texttt{lhs\_design\_matrix.npz} \\
\addlinespace
1 & Seed bank creation &
  Seeds config YAML &
  \texttt{seeds/nucleon-seeds.hdf} \\
\addlinespace
2 & Nuclear configuration &
  \texttt{nucleon-seeds.hdf}, \newline isobar config YAML &
  \texttt{design\_NNNN/WS1.hdf}, \newline \texttt{design\_NNNN/WS2.hdf} \\
\addlinespace
3 & \trento simulation &
  \texttt{WS*.hdf}, fixed Trento params &
  \texttt{design\_NNNN/trento\_events.npy} \\
\addlinespace
4 & Eccentricity computation &
  \texttt{trento\_events.npy} &
  \texttt{results/observables.npz} \\
\addlinespace
5 & Visualisation &
  \texttt{observables.npz} &
  \texttt{results/plots/*.pdf} \\
\bottomrule
\end{tabular}
\end{table}

\subsection{Stage 0: Latin Hypercube Generation}
\label{subsec:stage0}

The implementation of the maximin LHS follows
\cref{alg:lhs_generation}. For $N_\mathrm{d} = 100$, $d = 5$, and
$N_\mathrm{iter} = 2000$, the algorithm completes in
$\mathcal{O}(1)\,\si{s}$ on a modern CPU.

\begin{algorithm}[H]
\caption{Maximin-Optimised Latin Hypercube Sampling}
\label{alg:lhs_generation}
\SetAlgoLined
\KwIn{$N$ (design points), $d$ (parameters), seed, $N_\mathrm{iter}$ (trials)}
\KwOut{$\mathbf{U}^* \in [0,1]^{N \times d}$ (unit-cube design)}
$d_\mathrm{best} \leftarrow -\infty$\;
\For{$t = 1$ \KwTo $N_\mathrm{iter}$}{
  \ForEach{column $j = 1,\ldots,d$}{
    Draw $\pi_j \leftarrow \mathrm{Permutation}(N)$\;
    $U_{ij} \leftarrow (\pi_j[i] + u_{ij}) / N$,\;
    \hspace{1cm} where $u_{ij} \sim \mathcal{U}[0,1]$
  }
  $d_\mathrm{min} \leftarrow \min_{i \neq j} \|\bm{u}_i - \bm{u}_j\|_2$\;
  \If{$d_\mathrm{min} > d_\mathrm{best}$}{
    $d_\mathrm{best} \leftarrow d_\mathrm{min}$;\;
    $\mathbf{U}^* \leftarrow \mathbf{U}$\;
  }
}
\Return $\mathbf{U}^*$
\end{algorithm}

\subsection{Stages 1--2: Nuclear Configuration Generation}
\label{subsec:stages12}

For each design point $\bm{\theta}^{(i)}$, the following two-step
procedure generates the nuclear HDF5 files for \trento:

\paragraph{Step 1 (once per campaign):}
The seeds configuration is written as:
\begin{verbatim}
isobar_seeds:
  number_nucleons:  {value: 208}
  number_configs:   {value: 10000}
  output_file:      {filename: nucleon-seeds.hdf}
  number_of_parallel_processes: {value: -1}
\end{verbatim}
and \texttt{make\_seeds.py} is invoked.

\paragraph{Step 2 (once per design point):}
The isobar configuration file maps the LHS sample to the
isobar-sampler YAML schema.  Both isobars (\texttt{WS1},
\texttt{WS2}) are assigned identical parameters, reflecting the
symmetry of \PbPb collisions:
\begin{verbatim}
isobar_properties:
  isobar1:
    isobar_name: WS1
    WS_radius:            {value: R[i]}
    WS_diffusiveness:     {value: a[i]}
    beta_2:               {value: 0.0}
    gamma:                {value: 0.0}
    beta_3:               {value: beta3[i]}
    beta_4:               {value: beta4[i]}
    correlation_length:   {value: 0.4}
    correlation_strength: {value: -1.0}
  isobar2:
    isobar_name: WS2
    [identical parameters]
\end{verbatim}

\subsection{Stage 3: \trento Simulation}
\label{subsec:stage3}

\trento is called for each design point with the flag
\texttt{-{}-random-seed $= 1000 + i$} to ensure independent event
sequences across design points.  The \texttt{-o} flag is deliberately
\emph{omitted} so that \trento writes only per-event summary lines to
\texttt{stdout} rather than individual event files, eliminating
$N_\mathrm{ev}$ costly disk writes.  The \texttt{stdout} format is:
\begin{verbatim}
  event_id   b   N_part   M   epsilon_2   epsilon_3   epsilon_4   epsilon_5
\end{verbatim}
and the parsed array of shape
$(N_\mathrm{ev}, 8)$ is cached as a compressed \texttt{.npy} file.

\subsection{Stage 4: Eccentricity Computation and Bootstrapping}
\label{subsec:stage4}

After centrality assignment via \cref{eq:centrality}, the
$N_\mathrm{UC}$ ultracentral events are selected.  All observables
in \cref{tab:observables} are computed analytically from the sample
moments $m_k \equiv N_\mathrm{UC}^{-1}\sum_{i=1}^{N_\mathrm{UC}}
\varepsilon_2^k$.

Statistical uncertainties are estimated by the bootstrap
\cite{Efron1979}: $N_\mathrm{boot} = 300$ resamples with replacement
are drawn from the $N_\mathrm{UC}$ events, observables are recomputed
for each resample, and the $16^\mathrm{th}$--$84^\mathrm{th}$
percentile interval of the resulting bootstrap distribution is
reported as the $\pm 1\sigma$ band.  This correctly propagates
the wider uncertainty in the small $0$--$0.1\%$ bin relative to the
$0$--$1\%$ bin.

\begin{algorithm}[H]
\caption{Bootstrap Uncertainty Estimation for Ultracentral Observables}
\label{alg:bootstrap}
\SetAlgoLined
\KwIn{Events $\{e_1,\ldots,e_{N_\mathrm{UC}}\}$,\; $N_\mathrm{boot}$}
\KwOut{Central values $\hat{y}$, lower $\hat{y}_-$, upper $\hat{y}_+$}
$\hat{y} \leftarrow \textsc{ComputeObs}(\{e_i\})$\;
\For{$b = 1$ \KwTo $N_\mathrm{boot}$}{
  $\mathcal{I} \leftarrow$ sample $N_\mathrm{UC}$ indices with replacement\;
  $\hat{y}^{(b)} \leftarrow \textsc{ComputeObs}(\{e_{\mathcal{I}}\})$\;
}
$\hat{y}_- \leftarrow \mathrm{Percentile}_{16}\bigl(\{\hat{y}^{(b)}\}\bigr)$\;
$\hat{y}_+ \leftarrow \mathrm{Percentile}_{84}\bigl(\{\hat{y}^{(b)}\}\bigr)$\;
\Return $(\hat{y},\, \hat{y}_-,\, \hat{y}_+)$
\end{algorithm}

% ──────────────────────────────────────────────────────────────────────────────
\section{Parameter Space and Prior Ranges}
\label{sec:priors}
% ──────────────────────────────────────────────────────────────────────────────

\Cref{tab:prior_ranges} summarises the five free parameters, their
prior ranges, physical motivation, and key references.

\begin{table}[htbp]
\centering
\caption{Free parameters, prior ranges, central (reference) values,
  and literature references for the \PbPb ultracentral scan.}
\label{tab:prior_ranges}
\renewcommand{\arraystretch}{1.25}
\begin{tabular}{lccccl}
\toprule
Parameter & Symbol & Range & Central & Unit & Reference \\
\midrule
WS radius     & $R$      & $[6.50,\,6.80]$ & $6.624$ & fm & \cite{deVries1987} \\
WS diffuseness & $a$     & $[0.44,\,0.65]$ & $0.549$ & fm & \cite{deVries1987,PREX2021} \\
Octupole      & $\beta_3$ & $[0.00,\,0.12]$ & $0.0$ & — & \cite{Carzon2020,Xu2025} \\
Hexadecapole  & $\beta_4$ & $[-0.02,\,0.06]$ & $0.0$ & — & \cite{Bally2022} \\
Nucleon width & $w$      & $[0.40,\,1.20]$ & $0.76$ & fm & \cite{Bernhard2019} \\
\bottomrule
\end{tabular}
\end{table}

\noindent
The octupole range $\beta_3 \in [0, 0.12]$ is chosen to span the
constraint of \cite{Carzon2020} ($\beta_3 \lesssim 0.0375$, which
disfavours a static octupole as the sole resolution of the puzzle)
and the more liberal range entertained by proposals involving a
dynamical octupole ``breathing mode'' \cite{Xu2025}.  The upper end
of the hexadecapole range $\beta_4 = 0.06$ exceeds current DFT
predictions but is included to assess the sensitivity of $v_4\{4\}^4$
to large $\beta_4$ values.

% ──────────────────────────────────────────────────────────────────────────────
\section{Implementation Details and Code Structure}
\label{sec:implementation}
% ──────────────────────────────────────────────────────────────────────────────

\subsection{Directory Structure}
\label{subsec:dir_structure}

The pipeline creates the following directory tree under
\texttt{pbpb\_scan\_276GeV/} (or \texttt{...502GeV/}
for $\sqrtsnn = 5.02\,\si{TeV}$):

\begin{verbatim}
pbpb_scan_276GeV/
  lhs_design_matrix.npz          # design matrix + metadata
  seeds/
    seeds-conf.yaml               # isobar-sampler seed config
    nucleon-seeds.hdf             # seed bank (shared)
  design_0000/
    isobar-conf.yaml              # per-design-point isobar config
    WS1.hdf, WS2.hdf              # nuclear configurations
    trento_events.npy             # cached Trento output
  design_0001/ ...
  results/
    observables.npz               # all observables × all design pts
    summary_table.txt             # human-readable table
    plots/
      corner_lhs.pdf              # LHS space-filling check
      obs_all_panels.pdf          # full observable grid
      obs_vs_R.pdf, ...           # per-parameter scatter plots
\end{verbatim}

\subsection{Key Code Sections}
\label{subsec:code_sections}

The eccentricity cumulant computation (\texttt{stage4\_eccentricities})
implements \cref{eq:e2_cumulant,eq:e4_cumulant,eq:e6_cumulant,eq:e8_cumulant}
as vectorised NumPy operations.  The following listing shows the core
moment computation and cumulant evaluation:

\begin{lstlisting}[caption={Core eccentricity cumulant computation
  (from \texttt{pbpb\_ultracentral\_scan.py}).},
  label=lst:cumulants]
def _moments(eps, orders=(2, 4, 6, 8)):
    """Compute <eps^n> for each n in a single pass."""
    return {n: np.mean(eps**n) for n in orders}

def _e2_cumulants(m):
    """Return (e{2}, e{4}, e{6}, e{8}) from moment dict m."""
    m2, m4, m6, m8 = m[2], m[4], m[6], m[8]

    e2 = np.sqrt(m2)

    inner4 = 2.0 * m2**2 - m4          # e{4}^4 = 2<e^2>^2 - <e^4>
    e4 = inner4**0.25 if inner4 >= 0 else np.nan

    c6 = m6 - 9.0*m4*m2 + 12.0*m2**3  # 6-particle cumulant
    e6 = (-c6/4.0)**(1/6) if -c6/4 >= 0 else np.nan

    c8 = (m8 - 16.0*m6*m2 - 18.0*m4**2
          + 144.0*m4*m2**2 - 144.0*m2**4)
    e8 = (c8/33.0)**(1/8) if c8/33 >= 0 else np.nan

    return e2, e4, e6, e8
\end{lstlisting}

The hexadecapole sign-change observable and nonlinear coupling
(\cref{eq:e4_4c,eq:sign_change_trigger}) are computed as:

\begin{lstlisting}[caption={Hexadecapole cumulant and sign-change trigger.},
  label=lst:puzzle_obs]
# e4{4}^4 = 2<e4^2>^2 - <e4^4>  (negative in ultracentral: the puzzle)
m4_e4  = np.mean(eps4**4)
e4_4c  = float(2 * np.mean(eps4**2)**2 - m4_e4)

# Sign-change trigger: > 2 => nonlinear channel drives v4{4}^4 < 0
nl_r8  = float(m[8] / m[4]**2) if m[4] > 0 else np.nan
\end{lstlisting}

\subsection{Configuration Flags}
\label{subsec:config_flags}

Before running, the user must set:
\begin{itemize}[leftmargin=*]
  \item \texttt{SQRTS}: collision energy (2.76 or 5.02, in TeV);
  \item \texttt{ISOBAR\_SAMPLER\_DIR}: path to the local clone of
        \texttt{mluzum/Isobar-Sampler};
  \item \texttt{TRENTO\_BIN}: path to the compiled \trento binary;
  \item \texttt{N\_DESIGN}: number of design points (default 100);
  \item \texttt{N\_EVENTS}: Trento events per design point
        (default $5\times 10^5$).
\end{itemize}
Each stage can be skipped independently with the
\texttt{RUN\_STAGE\_*} boolean flags, enabling re-analysis of
previously computed events without re-running \trento.

% ──────────────────────────────────────────────────────────────────────────────
\section{Summary}
\label{sec:summary}
% ──────────────────────────────────────────────────────────────────────────────

This chapter has presented the theoretical and computational
foundations of a systematic scan of the \PbPb ultracentral flow
puzzle.  The principal elements are:

\begin{enumerate}[leftmargin=*]
  \item A deformed Woods-Saxon model of the $^{208}$Pb nucleus with
        five free parameters ($R$, $a$, $\beta_3$, $\beta_4$, $w$)
        spanning literature-motivated prior ranges
        (\cref{tab:prior_ranges}).
  \item The isobar-sampler seed-bank method of \cite{Luzum2023_sampler},
        which ensures that configurations generated at different design
        points share the same underlying nucleon-position fluctuations,
        isolating the genuine parameter response from statistical noise.
  \item The \trento entropy-deposition model \cite{Moreland2015} with
        fixed Bayesian-calibrated parameters \cite{Bernhard2019,Nijs2021}.
  \item A 100-point maximin Latin Hypercube design in five dimensions,
        ensuring uniform marginal coverage and good pairwise separation
        (\cref{alg:lhs_generation}).
  \item A complete set of eleven eccentricity observables including
        the four puzzle-critical quantities $T_{32}$, $R_{42}$,
        $\Gamma_2$, and $\varepsilon_4\{4\}^4$ (\cref{tab:observables}).
  \item Bootstrap uncertainty estimation with $N_\mathrm{boot} = 300$
        resamples providing honest $\pm 1\sigma$ bands on all
        observables (\cref{alg:bootstrap}).
\end{enumerate}

The output design matrix, observables, and their uncertainties will
serve as training data for a Gaussian Process emulator in the
following chapter, enabling a full Bayesian posterior analysis of
the puzzle parameters.

% ──────────────────────────────────────────────────────────────────────────────
\section*{Acknowledgements}
% ──────────────────────────────────────────────────────────────────────────────

The author thanks M.~Luzum (USP) for providing the Isobar-Sampler
code and for illuminating discussions on short-range correlations in
nucleus-nucleus collisions.

% ══════════════════════════════════════════════════════════════════════════════
%% BIBLIOGRAPHY (BibTeX entries — replace with your .bib file in thesis)
% ══════════════════════════════════════════════════════════════════════════════
\begin{thebibliography}{99}

\bibitem{Aamodt2011}
  K.~Aamodt \textit{et al.} (ALICE Collaboration),
  Phys.\ Rev.\ Lett.\ \textbf{107}, 032301 (2011)
  [arXiv:1105.3865].

\bibitem{Aad2012}
  G.~Aad \textit{et al.} (ATLAS Collaboration),
  Phys.\ Rev.\ C \textbf{86}, 014907 (2012)
  [arXiv:1203.3087].

\bibitem{Luzum2013}
  M.~Luzum and J.-Y.~Ollitrault,
  Phys.\ Rev.\ Lett.\ \textbf{111}, 132301 (2013)
  [arXiv:1305.4894].

\bibitem{Adam2016}
  J.~Adam \textit{et al.} (ALICE Collaboration),
  Phys.\ Rev.\ Lett.\ \textbf{117}, 182301 (2016)
  [arXiv:1604.07663].

\bibitem{Aaboud2019}
  M.~Aaboud \textit{et al.} (ATLAS Collaboration),
  Phys.\ Rev.\ C \textbf{100}, 034406 (2019)
  [arXiv:1901.07808].

\bibitem{Sirunyan2017}
  A.~M.~Sirunyan \textit{et al.} (CMS Collaboration),
  Phys.\ Lett.\ B \textbf{789}, 643 (2019)
  [arXiv:1711.05594].

\bibitem{Giacalone2017}
  G.~Giacalone, J.~Noronha-Hostler and J.-Y.~Ollitrault,
  Phys.\ Rev.\ C \textbf{95}, 054910 (2017)
  [arXiv:1702.01730].

\bibitem{Chatrchyan2014}
  S.~Chatrchyan \textit{et al.} (CMS Collaboration),
  Phys.\ Rev.\ C \textbf{89}, 044906 (2014)
  [arXiv:1310.8651].

\bibitem{Carzon2020}
  P.~Carzon, S.~Rao, M.~Luzum, M.~Sievert and J.~Noronha-Hostler,
  Phys.\ Rev.\ C \textbf{102}, 054905 (2020)
  [arXiv:2007.00780].

\bibitem{Giacalone2017b}
  G.~Giacalone, J.~Noronha-Hostler and J.-Y.~Ollitrault,
  Phys.\ Rev.\ C \textbf{96}, 054903 (2017)
  [arXiv:1707.08765].

\bibitem{Luzum2023}
  M.~Luzum \textit{et al.},
  Quark Matter 2023 proceedings,
  arXiv:2312.10129 (2023).

\bibitem{WoodsSaxon1954}
  R.~D.~Woods and D.~S.~Saxon,
  Phys.\ Rev.\ \textbf{95}, 577 (1954).

\bibitem{deVries1987}
  H.~de~Vries, C.~W.~de~Jager and C.~de~Vries,
  At.\ Data Nucl.\ Data Tables \textbf{36}, 495 (1987).

\bibitem{PREX2021}
  D.~Adhikari \textit{et al.} (PREX Collaboration),
  Phys.\ Rev.\ Lett.\ \textbf{126}, 172502 (2021)
  [arXiv:2102.10767].

\bibitem{Roca-Maza2018}
  X.~Roca-Maza and N.~Paar,
  Prog.\ Part.\ Nucl.\ Phys.\ \textbf{101}, 96 (2018).

\bibitem{Bally2022}
  B.~Bally \textit{et al.},
  Phys.\ Rev.\ Lett.\ \textbf{128}, 082301 (2022)
  [arXiv:2108.09578].

\bibitem{Xu2025}
  B.~Xu \textit{et al.},
  arXiv:2504.19644 (2025).

\bibitem{Teaney2011}
  D.~Teaney and L.~Yan,
  Phys.\ Rev.\ C \textbf{83}, 064904 (2011)
  [arXiv:1010.1876].

\bibitem{Borghini2001}
  N.~Borghini, P.~M.~Dinh and J.-Y.~Ollitrault,
  Phys.\ Rev.\ C \textbf{64}, 054901 (2001)
  [arXiv:nucl-ex/0007007].

\bibitem{Moreland2015}
  J.~S.~Moreland, J.~E.~Bernhard and S.~A.~Bass,
  Phys.\ Rev.\ C \textbf{92}, 011901 (2015)
  [arXiv:1412.4708].

\bibitem{Bernhard2019}
  J.~E.~Bernhard, J.~S.~Moreland and S.~A.~Bass,
  Nat.\ Phys.\ \textbf{15}, 1113 (2019)
  [arXiv:1901.07808].

\bibitem{Nijs2021}
  G.~Nijs, W.~van der Schee, U.~G\"ursoy and R.~Snellings,
  Phys.\ Rev.\ C \textbf{103}, 054909 (2021)
  [arXiv:2010.15134].

\bibitem{ALICE_2.76}
  K.~Aamodt \textit{et al.} (ALICE Collaboration),
  Phys.\ Rev.\ Lett.\ \textbf{105}, 252301 (2010)
  [arXiv:1011.3916].

\bibitem{ALICE_5.02}
  J.~Adam \textit{et al.} (ALICE Collaboration),
  Phys.\ Lett.\ B \textbf{772}, 567 (2017)
  [arXiv:1612.08966].

\bibitem{Luzum2023_sampler}
  M.~Luzum \textit{et al.},
  Eur.\ Phys.\ J.\ A \textbf{59}, 110 (2023)
  [arXiv:2302.14026].

\bibitem{Morris1995}
  M.~D.~Morris and T.~J.~Mitchell,
  J.\ Stat.\ Plan.\ Inference \textbf{43}, 381 (1995).

\bibitem{Loeppky2009}
  J.~L.~Loeppky, J.~Sacks and W.~J.~Welch,
  Technometrics \textbf{51}, 366 (2009).

\bibitem{Efron1979}
  B.~Efron,
  Ann.\ Stat.\ \textbf{7}, 1 (1979).

\end{thebibliography}

\end{document}
